% Taxonomy of data-driven flare forecasting; where this thesis sits.
\begin{figure}[!htbp]
\centering
\begin{tikzpicture}[
  font=\small, >=Latex, node distance=6mm and 9mm,
  card/.style={draw, rounded corners=3pt, minimum height=13mm, text width=30mm,
               align=center, fill=black!4, inner sep=3pt},
  mine/.style={card, fill=green!14, very thick},
]
\node[card] (sharp) {\textbf{ML on SHARP scalars}\\\scriptsize SVM on integrated
  field parameters~\cite{bobra2015}};
\node[card, right=of sharp] (cnn) {\textbf{Deep spatial}\\\scriptsize CNN on
  magnetogram images~\cite{huang2018}};
\node[card, right=of cnn] (lstm) {\textbf{Deep temporal}\\\scriptsize LSTM on SHARP
  time series~\cite{liu2019,nishizuka2018}};
\node[card, below=14mm of sharp] (phys) {\textbf{Winding physics}\\\scriptsize signed
  topological precursor~\cite{mactaggart2021}};
\node[card, below=14mm of cnn] (scalar) {\textbf{Winding $\to$ scalars}\\\scriptsize
  integral + XGBoost~\cite{williams2026}};
\node[mine, below=14mm of lstm] (mine) {\textbf{Winding \emph{map} (this thesis)}\\
  \scriptsize 2D structure, spatiotemporal model};
\draw[->] (sharp) -- (cnn); \draw[->] (cnn) -- (lstm);
\draw[->] (phys) -- (scalar);
\draw[->, very thick, green!45!black] (scalar) -- (mine);
\draw[->, dashed] (lstm) -- (mine);
\draw[->, dashed] (phys) to[out=90,in=180] (sharp);
\end{tikzpicture}
\caption[Forecasting landscape]{\textbf{Where this thesis sits.} Data-driven
flare forecasting progressed from machine learning on integrated SHARP scalars
to deep spatial and deep temporal models. In parallel, magnetic winding emerged
as a physically motivated precursor, but existing work reduces the winding
\emph{map} to a single spatial integral, in which oppositely signed regions
cancel. This thesis instead feeds the two-dimensional winding map, with its
structure intact, to a spatiotemporal model.}
\label{fig:taxonomy}
\end{figure}
